\subsection{有限表示模}

设 \(\mathbf{A}\) 是环。左（相应地右）A-模的一个正合列

\[
\mathrm{L} _ {1} \rightarrow \mathrm{L} _ {0} \rightarrow \mathrm{E} \rightarrow 0\tag{6}
\]

其中 \(L_{0}\) 和 \(L_{1}\) 是自由的，称为左（相应地右）A-模 E 的一个表示（或长度为 1 的表示）。

每个 A-模 E 都有一个表示。事实上我们知道（《代数学》第二章 §1 第 11 目命题 20），存在一个满同态 \(u \colon \mathbf{L}_0 \to \mathbf{E}\) ，其中 \(\mathbf{L}_0\) 是自由的；若 \(\mathbf{R}\) 是 \(u\) 的核，则类似地存在满同态 \(v \colon \mathbf{L}_1 \to \mathbf{R}\) ，其中 \(\mathbf{L}_1\) 是自由的。若把 \(v\) 视为 \(\mathbf{L}_1\) 到 \(\mathbf{L}_0\) 的同态，则序列 \(\mathbf{L}_1 \xrightarrow{v} \mathbf{L}_0 \xrightarrow{u} 0\) 按定义是正合的，由此即得我们的断言。

若 \(\rho: \mathbf{A} \to \mathbf{B}\) 是环同态，则 \(\mathbf{E}\) 的每个表示 (6) 都诱导 \(\mathbf{E}_{(\mathbf{B})} = \mathbf{E} \otimes_{\mathbf{A}} \mathbf{B}\) 的一个表示：

\[
\mathrm{L} _ {1} \otimes_ {\mathrm{A}} \mathrm{B} \rightarrow \mathrm{L} _ {0} \otimes_ {\mathrm{A}} \mathrm{B} \rightarrow \mathrm{E} \otimes_ {\mathrm{A}} \mathrm{B} \rightarrow 0\tag{7}
\]

这由第 1 目的引理 1 以及 \(\mathbf{L} \otimes_{\mathbf{A}} \mathbf{B}\) 在 \(\mathbf{L}\) 自由时是自由 B-模这一事实得出。

模 E 的表示 (6) 称为有限的，如果自由模 \(L_{0}\) 和 \(L_{1}\) 都有有限基。显然，若表示 (6) 有限，则表示 (7) 也有限。称 E 为有限表示 A-模，如果它有一个有限表示。

断言 (i) 由定义直接得出。若 A 是左 Noether 环，且存在满同态 \(u \colon \mathbf{L}_0 \to \mathbf{E}\) ，其中 \(\mathbf{L}_0\) 是有有限基的自由左 A-模，则 \(u\) 的核 R 是有限生成的（《代数学》第八章 §2 第 1 目命题 1 及第 3 目命题 7），因而存在满同态 \(v \colon \mathbf{L}_1 \to \mathbf{R}\) ，其中 \(\mathbf{L}_1\) 是自由的且有有限基，而正合列 \(\mathbf{L}_1 \xrightarrow{v} \mathbf{L}_0 \xrightarrow{u} \mathbf{E} \to 0\) 是 E 的一个有限表示；由此得 (ii)。

最后，设 E 是一个有限生成的投射模；则它是有限生成自由模 \(L_{0}\) 的一个直因子（《代数学》第二章 §2 第 2 目命题 4 的推论）；满同态 \(L_{0} \rightarrow E\) 的核 R 于是同构于 \(L_{0}\) 的一个商，因而是有限生成的，然后像 (ii) 那样完成证明。

引理 9. 设 A 是环，E 是有限表示 A-模。对每个正合列

\[
0 \longrightarrow \mathrm{F} \stackrel {j} {\longrightarrow} \mathrm{G} \stackrel {p} {\longrightarrow} \mathrm{E} \longrightarrow 0
\]

其中 G 是有限生成的，模 F 是有限生成的。

设 \(L_{1} \stackrel{r}{\rightarrow} L_{0} \stackrel{s}{\rightarrow} E \rightarrow 0\) 是一个有限表示；若 \((e_{i})\) 是 \(L_{0}\) 的一个基，则对每个 i 存在元素 \(g_{i} \in G\) 使得 \(p(g_{i}) = s(e_{i})\) ；于是同态 \(u: L_{0} \rightarrow G\) （它对所有 i 满足 \(u(e_{i}) = g_{i}\) ）满足 \(s = p \circ u\) 。由于 \(s \circ r = 0\) ，我们有 \(u(r(L_{1})) \subset \operatorname{Ker} p\) ，且由于 \(\operatorname{Ker} p\) 同构于 F，可以看出存在同态 \(v: L_{1} \rightarrow F\) 使得图

\[
\begin{array}{c} \mathrm{L} _ {1} \xrightarrow {r} \mathrm{L} _ {0} \xrightarrow {s} \mathrm{E} \longrightarrow 0 \\ v \Biggl \downarrow \qquad \qquad u \Biggl \downarrow \qquad \qquad 1 _ {\mathrm{E}} \Biggl \downarrow \\ \mathrm{F} \xrightarrow {} _ {j} \mathrm{G} \xrightarrow {} _ {p} \mathrm{E} \longrightarrow 0 \end{array}
\]

是交换的。由于 \(j\) 是单射而 \(s\) 是满射，我们可以应用蛇形图（§1 第 4 目，命题 4），换句话说，存在正合列

\[
0 = \operatorname{Ker} 1 _ {\mathrm{E}} \xrightarrow {d} \text { Coker } v \longrightarrow \text { Coker } u \longrightarrow \text { Coker } 1 _ {\mathrm{E}} = 0.
\]

这表明 Coker v 同构于 \(G/u(L_{0})\) ，后者由假设是有限生成的。此外我们有正合列

\[
0 \rightarrow v (\mathrm{L} _ {1}) \rightarrow \mathrm{F} \rightarrow \text { Coker } v \rightarrow 0
\]

且由于 \(v(\mathbf{L}_1)\) 和 Coker \(v\) 都是有限生成的，F 也是有限生成的（《代数学》第二章 §1 第 7 目命题 9 的推论 5）。

